\section{Software}

The implementation uses PyTorch and Hydra. The canonical workflow consists of balanced-dataset materialization and audit, sharded token-level corpus collection, independent actor and critic CLT training, global replacement evaluation, target-specific graph construction, and original-model intervention sweeps. Every corpus shard is tensor-only and hash-verified; accompanying JSONL metadata preserves source provenance and source-group-disjoint splits. Model specifications record the source-checkpoint hash, corpus-manifest hash, natural residual-stream preprocessing contract, CLT architecture, and branch identity. On Jean Zay, all generated data, checkpoints, metrics, graphs, logs, caches, and temporary files are written under the project \texttt{SCRATCH} directory and remain separate from committed configs and conclusions on \texttt{WORK}.

The initial CLT configuration uses 1,024 features at each of eight path layers, initial JumpReLU threshold $0.01$, surrogate-gradient bandwidth $10^{-3}$, and decoder-bundle normalization after every optimizer step. Each branch is trained for at most 100,000 AdamW steps with batch size 4,096, learning rate $3\times10^{-4}$, sparsity coefficient $10^{-3}$, unit tanh scale, and gradient-norm clipping at 1.0. Corpus tensors are stored in float16 and converted to the model parameter dtype for optimization and evaluation. Validation is logged every 500 steps on at most 100,000 token rows, with checkpoints every 5,000 steps.

At 16,588,800 token rows, the six-tensor corpus schema occupies approximately 152.9 GB (142.4 GiB) before JSONL metadata and filesystem overhead. Corpus tensors and JSONL metadata are written to Jean Zay's SCRATCH, while a small copy of the corpus and run manifests, trained CLTs, metrics, and graphs remains on WORK.

Global replacement is evaluated on 100 pooled held-out batches and reported overall and per environment. Prompt-local graphs retain 80\% of indirect feature-node influence followed by 98\% of candidate edge influence. The intervention pilot applies factors $\{-2,-1,-0.5,0.5,1,2\}$ to the selected feature's complete learned write bundle and compares eight seeded random bundles with matched per-target decoder norms; controls are paired post hoc by nearest induced policy KL. These values define the initial engineering pilot and may be revised from training and validation health only, before the held-out test analysis.
